```latex
\documentclass[11pt]{amsart}
\usepackage{amssymb,amsmath,amsthm,amscd}

\newtheorem{theorem}{Theorem}[section]
\newtheorem{lemma}[theorem]{Lemma}
\newtheorem{proposition}[theorem]{Proposition}
\newtheorem{corollary}[theorem]{Corollary}
\newtheorem{definition}[theorem]{Definition}
\newtheorem{remark}[theorem]{Remark}
\newtheorem{example}[theorem]{Example}

\newcommand{\Sph}{\mathbb{S}^{2}}
\newcommand{\Cpt}{\widehat{\mathbb{C}}}
\newcommand{\Teich}{\operatorname{Teich}}
\newcommand{\Mod}{\operatorname{Mod}}
\newcommand{\PMod}{\operatorname{PMod}}
\newcommand{\Def}{\mathcal{D}}
\newcommand{\cA}{\mathbf{c}}
\newcommand{\Homeo}{\operatorname{Homeo}}

\begin{document}

\title{A Dehn twist acting with infinite orbits on the components\\
of a deformation space of a quadratic rational map}

\author{}
\date{}

\begin{abstract}
Let $f:\Sph\to\Sph$ be a quadratic rational map whose critical point $c_1$
lies on a cycle of length $4$, and whose second critical value $v=f(c_2)$
does not meet the cycle. With $A$ the $4$-cycle, $B=A\cup\{v\}$, we consider
the deformation space $\Def_f=\{\tau\in\Teich(\Sph,B)\;:\; f^*\tau=i_*\tau\}$.
We describe an explicit simple closed curve
\[
\gamma \;=\; \partial D,\qquad D\cap B=\{c_1,\;f^{\circ 2}(c_1)\},
\]
(i.e.\ the curve enclosing $c_1$ and $f^{\circ 2}(c_1)$, or equivalently the
complementary curve enclosing $v$, $f(c_1)$ and $f^{\circ 3}(c_1)$), such
that the Dehn twist $T_\gamma$ preserves $\Def_f$, and such that the cyclic
group $\langle T_\gamma\rangle$ acts on the set $\pi_0(\Def_f)$ of connected
components of $\Def_f$ with infinite orbits. In particular $\Def_f$ has
infinitely many connected components, recovering the main theorem of
\cite{Luo}.
\end{abstract}

\maketitle

\section{Setup and statement of the main theorem}

\subsection{Deformation spaces}
Let $f:\Sph\to\Sph$ be an orientation-preserving branched cover, let
$V(f)$ denote its set of critical values, and let $A,B\subset\Sph$ be finite
sets with $A\cup f(A)\cup V(f)\subset B$. Following Epstein
\cite{Epstein} we define
\[
\sigma_f:=f^*:\Teich(\Sph,B)\longrightarrow \Teich(\Sph,A),
\qquad
\sigma^B_A:=i_*:\Teich(\Sph,B)\longrightarrow \Teich(\Sph,A),
\]
where $f^*$ is obtained by pulling back a complex structure of
$(\Sph,B)$ through the covering $f:(\Sph,f^{-1}(B))\to(\Sph,B)$ (this is
well defined because $V(f)\subset B$) and then erasing the points of
$f^{-1}(B)\setminus A$, and where $i_*$ forgets the points of $B\setminus A$.
The \emph{deformation space} of $f$ is the equalizer
\[
\Def_f \;:=\;\{\tau\in\Teich(\Sph,B)\;:\;\sigma_f(\tau)=\sigma^B_A(\tau)\}.
\]
Throughout, $\Teich(\Sph,X)$ denotes the Teichm\"uller space of the sphere
with the points of $X$ marked \emph{pointwise}, so that
$\Teich(\Sph,X)$ is a simply connected complex manifold of dimension
$|X|-3$, and the \emph{pure} mapping class group
$\PMod(\Sph,X)=\pi_0\Homeo^+(\Sph,X)$ (isotopy classes of homeomorphisms
fixing $X$ pointwise) acts freely and holomorphically on the right on
$\Teich(\Sph,X)$, with quotient the (ordered) configuration moduli space
$\mathcal{M}_{0,|X|}$. We write the right action as $\tau\mapsto\tau\cdot h$.

The deformation space is nonempty exactly when $f$ is equivalent (in the
sense of Thurston, rel.\ $A$) to a rational map \cite{Epstein,Luo}; in that
case $\Def_f$ parametrizes marked rational maps combinatorially equivalent
to $f$.

\subsection{The family}
Let $f$ be a \emph{quadratic rational map} with critical points $c_1,c_2$.
Assume:
\begin{enumerate}
\item[(H1)] $c_1$ is periodic of period $4$: writing
$p_1=f(c_1),\ p_2=f^{\circ 2}(c_1),\ p_3=f^{\circ 3}(c_1)$, the points
$c_1,p_1,p_2,p_3$ are distinct and $f^{\circ 4}(c_1)=c_1$;
\item[(H2)] the second critical value $v:=f(c_2)$ is \emph{not} on the
cycle: $v\notin A:=\{c_1,p_1,p_2,p_3\}$.
\end{enumerate}
Thus $\langle f\rangle$ lies in the locus $\operatorname{Per}_4(0)^*\subset
\mathcal{M}_2$ of conjugacy classes of quadratic rational maps with a
superattracting $4$-cycle containing exactly one critical point
\cite{Luo}. We set
\[
A=\{c_1,p_1,p_2,p_3\},\qquad B=A\cup\{v\},\qquad \Def_f=\Def(f,A,B).
\]
Since $|B|-|A|=1$, Epstein's local theorem (see \cite{Epstein},
\cite[Thm.~1.2]{Luo}) gives:

\begin{theorem}[{Epstein; \cite[Thm.~1.2]{Luo}}]\label{T:dim}
For $f$ as above, $\Def_f\subset\Teich(\Sph,B)$ is a smooth complex
submanifold of dimension $|B|-|A|=1$; in particular $\Teich(\Sph,B)$ has
dimension $2$ and $\Teich(\Sph,A)\cong\mathbb{H}$ has dimension $1$.
\end{theorem}

\subsection{The curve and the main result}

\begin{definition}\label{D:curve}
Let $D\subset\Sph$ be a closed topological disk with
\[
D\cap B=\{c_1,\;p_2\}=\{c_1,\;f^{\circ 2}(c_1)\},
\qquad
\{p_1,p_3,v\}\subset \Sph\setminus D ,
\]
and let $\gamma:=\partial D$. Equivalently, $\gamma$ is the unique
isotopy class (rel.\ $B$) of essential simple closed curves in $(\Sph,B)$
which encloses $v$, $f(c_1)$ and $f^{\circ 3}(c_1)$ on one side and
$c_1$, $f^{\circ 2}(c_1)$ on the other. We let $T_\gamma$ denote the
(right-handed) Dehn twist about $\gamma$, an element of
$\PMod(\Sph,B)$.
\end{definition}

\begin{theorem}[Main Theorem]\label{T:main}
Let $f$ be a quadratic rational map satisfying {\rm (H1)}, {\rm (H2)}, let
$A,B$ and $\gamma$ be as above. Then:
\begin{enumerate}
\item $T_\gamma\in S_f$, where
$S_f:=\{h\in\PMod(\Sph,B)\;:\;\Def_f\cdot h=\Def_f\}$ is the stabilizer of
the deformation space. In particular the (first) power of the Dehn twist
about $\gamma$ preserves $\Def_f$.
\item For every $n\neq 0$ we have $T_\gamma^{\,n}\notin E_f$, where
$E_f:=\{h\in S_f\;:\;D_0\cdot h=D_0\}\subset S_f$ is the stabilizer (within
$S_f$) of a fixed connected component $D_0$ of $\Def_f$.
\item Consequently the $\langle T_\gamma\rangle$-orbit
$\{D_0\cdot T_\gamma^{\,n}\}_{n\in\mathbb{Z}}$ of any connected component
$D_0$ of $\Def_f$ is infinite, the index $[S_f:E_f]$ is infinite, and
$\Def_f$ has infinitely many connected components.
\end{enumerate}
\end{theorem}

Part (3) recovers \cite[Thm.~1.3]{Luo}. The proof combines a structure
theorem for $\Def_f$ (Sections~\ref{S:structure}--\ref{S:cosets}), a lifting
criterion for membership in $S_f$ (Section~\ref{S:lift}), the explicit
preimage computation for $\gamma$ (Section~\ref{S:curve}), and the index
computation of \cite{Luo} (Section~\ref{S:infinite}).

\section{Normalization of the family and the image of $\Def_f$}
\label{S:structure}

\subsection{The normalized model}
By Möbius conjugation we may normalize the marked $4$-cycle of a map in our
family as
\[
c_1\mapsto 0,\qquad p_1\mapsto 1,\qquad p_2\mapsto\infty,\qquad
p_3\mapsto t ,
\]
where $t\in\Cpt\setminus\{0,1,\infty\}$.

\begin{lemma}\label{L:family}
There is exactly one quadratic rational map $f_t$ with
$f_t(0)=1,\ f_t(1)=\infty,\ f_t(\infty)=t,\ f_t(t)=0$ and $0$ a critical
point (necessarily $f_t(z)-1$ vanishes doubly at $0$), namely
\[
f_t(z)\;=\;\frac{t\,z^2-(t^2+t+1)z+(t^2+t)}
{z^2-(t^2+t+1)z+(t^2+t)} .
\]
Its second critical point and second critical value are
\[
c_2(t)=\frac{2(t^2+t)}{t^2+t+1},
\qquad
v(t)=\frac{t^4-2t^3-t^2+2t+1}{t^4+2t^3-t^2-2t+1}.
\]
\end{lemma}

\begin{proof}
Write $f=N/D$ with $D$ monic of degree $2$. The conditions
$f(1)=\infty$, $f(\infty)=t$ and the doubling condition
$f(z)-1\equiv$ (const)$\cdot z^2$ at $0$ read $D(1)=0$, $\deg N-\deg D$
ratio $=t$, $N-D=(t-1)z^2$; solving $N(t)=0$ with $N(0)=D(0)$ then forces
the displayed normal form. The critical points of $f_t$ are the zeros of
$W(z)=N'(z)D(z)-N(z)D'(z)=(t-1)\,z\bigl(-(t^2+t+1)z+2(t^2+t)\bigr)$, i.e.\
$0$ and $c_2(t)$, and $v(t)=f_t(c_2(t))$ is a direct substitution
(simplifying with $s=t^2+t$, one finds $v(t)=\bigl((s+1)^2-4ts\bigr)/(s-1)^2$).
\end{proof}

Thus the (ordered) configuration moduli space is
\[
\mathcal{M}_{0,5}\;\cong\;\{(t,v)\;:\;t,v\in\Cpt\setminus\{0,1,\infty\},\ t\neq v\},
\]
with coordinates $(t,v)$ recording the positions of $p_3$ and of the free
critical value after the normalization $c_1\mapsto 0,\ p_1\mapsto 1,\
p_2\mapsto\infty$.

\subsection{The image curve}

Let
\[
C\;:=\;\{(t,v(t))\;:\; t\in \Cpt,\ v(t)\notin\{0,1,\infty\}\cup\{t\},\
t\notin\{0,1,\infty\}\}\;\subset\;\mathcal{M}_{0,5},
\]
the graph of the second critical value function, with the parameters
at which two marked points collide removed. By Lemma~\ref{L:family},
$C$ is the complement, in a rational curve, of the finite set
\[
P=\{0,1,\infty\}\cup Z\bigl(t^4-2t^3-t^2+2t+1\bigr)
\cup Z\bigl(t^4+2t^3-t^2-2t+1\bigr)\cup Z\bigl(t^5+t^4+t^3-t^2-3t-1\bigr),
\]
so $C\cong\Cpt\setminus P$ is a connected hyperbolic curve; $P$ records
precisely the degenerations $p_3\to c_1,p_1,p_2$ and
$v\to c_1,p_1,p_2,p_3$.

Let $\pi_B:\Teich(\Sph,B)\to\mathcal{M}_{0,5}$ be the universal covering
(with deck group $\PMod(\Sph,B)$), and $\pi_A:\Teich(\Sph,A)\to
\mathcal{M}_{0,4}\cong\Cpt\setminus\{0,1,\infty\}$ the universal covering
with deck group $\PMod(\Sph,A)$.

\begin{lemma}\label{L:image}
$\pi_B(\Def_f)=C$.
\end{lemma}

\begin{proof}
Fix $\tau=[\phi]\in\Def_f$, let $(t,v):=\pi_B(\tau)$, and normalize
$\phi(c_1)=0,\ \phi(p_1)=1,\ \phi(p_2)=\infty,\ \phi(p_3)=t,\
\phi(v)=v$. Put $g:=\phi\circ f$; then $Y:=(\Sph,g^*\text{std})$ is a
Riemann surface and $g:Y\to\Cpt$ is holomorphic of degree $2$, ramified at
$c_1$ (over $1$) and at $c_2$ (over $v$). Let $m:Y\to\Cpt$ be the
biholomorphism with $m(c_1)=0,\ m(p_1)=1,\ m(p_2)=\infty$, and put
$u:=m(p_3)$ and $\hat f:=g\circ m^{-1}$. Then $\hat f$ is a quadratic
rational map with $\hat f(0)=1$, $\hat f(1)=\infty$, $\hat f(\infty)=t$,
$\hat f(u)=0$, ramified at $0$ over $1$ and at
$\hat f^{-1}(v)$ over $v$. Since $\tau\in\Def_f$, the two
$\Teich(\Sph,A)$-points $\sigma_f(\tau)$ and $i_*\tau$ coincide, so in
particular their configurations agree: $u=t$. Then $\hat f(t)=0$ and $0$
is a critical point of $\hat f$ with $\hat f(0)=1$, so by
Lemma~\ref{L:family}, $\hat f=f_t$ and hence
$v=\hat f(\text{second critical point})=v(t)$. Thus
$\pi_B(\tau)\in C$. Conversely, if $t\in\Cpt\setminus P$ and $\beta:
\Sph\to\Cpt$ is any marking conjugating $f$ (in its Thurston class, rel.\
$B$) to the rational map $f_t$ with $\beta(B)=(0,1,\infty,t,v(t))$,
then the standard argument
(\cite{Epstein}; compare the proof of Lemma~\ref{L:liftcriterion} below)
shows $[\beta]\in\Def_f$; such markings exist because the family
$\{f_t\}_{t\in C}$ is topologically (indeed Thurston-)trivial over $C$.
Hence $\pi_B(\Def_f)=C$.
\end{proof}

\begin{lemma}\label{L:open}
$\Def_f$ is a union of connected components of
$\pi_B^{-1}(C)$.
\end{lemma}

\begin{proof}
By Theorem~\ref{T:dim} and Lemma~\ref{L:image}, $\Def_f$ is a
$1$-dimensional complex submanifold contained in the $1$-dimensional
complex submanifold $\pi_B^{-1}(C)$ (a union of covering surfaces over the
curve $C$). At any $\tau\in\Def_f$ the tangent line of $\Def_f$ lies in the
tangent line of $\pi_B^{-1}(C)$; since both are $1$-dimensional, the two
submanifolds coincide locally near $\tau$, so $\Def_f$ is open in
$\pi_B^{-1}(C)$. It is closed in $\pi_B^{-1}(C)$ because it is the
equalizer of two holomorphic maps into the Hausdorff space
$\Teich(\Sph,A)$. Being open and closed, it is a union of connected
components of $\pi_B^{-1}(C)$.
\end{proof}

\section{Components, cosets, and the groups $S_f$, $E_f$}\label{S:cosets}

Fix a point $c_0=(t_0,v(t_0))\in C$ and a lift $\tau_0\in\Def_f$ over it
(Lemma~\ref{L:image}). Let
\[
G:=\PMod(\Sph,B),\qquad
H:=j_*\bigl(\pi_1(C,c_0)\bigr)\subset G,
\]
where $j:C\hookrightarrow\mathcal{M}_{0,5}$ is the inclusion and
$j_*$ is computed through the identification
$\pi_1(\mathcal{M}_{0,5})\cong G$ given by point-pushing (Birman).

\begin{proposition}\label{P:cosets}
\begin{enumerate}
\item The connected components of $\pi_B^{-1}(C)$ correspond bijectively to
the right cosets $\{Hh:\;h\in G\}$, the component containing $\tau_0\cdot h$
corresponding to $Hh$; the right action of $G$ sends the coset $Hh$ to
$H(hh')$.
\item The connected components of $\Def_f$ are precisely the components
$Hh$ of $\pi_B^{-1}(C)$ which are contained in $\Def_f$; in particular
$H\subset S_f$, and
\[
E_f\;=\;\{h'\in S_f:\;Hh'=H\}\;=\;S_f\cap H .
\]
\item For $h\in S_f$ and a component $D=Hh$ of $\Def_f$, we have
$D\cdot h'=Dh'$; in particular the $\langle T_\gamma\rangle$-orbit of a
component $D_0$ is in bijection with the cosets
$\{HT_\gamma^{\,n}\}_{n\in\mathbb{Z}}$, and it is infinite if and only if
$T_\gamma^{\,n}\notin H$ for all $n\ne 0$.
\end{enumerate}
\end{proposition}

\begin{proof}
(1) Since $\pi_B$ is the universal covering with deck group $G$, and $C$
is connected and locally path connected, a point of
$\pi_B^{-1}(C)$ over $c_0$ has the form $\tau_0 h$ ($h\in G$), and two such
points $\tau_0h,\tau_0h'$ lie in the same component of $\pi_B^{-1}(C)$ if
and only if the projected loop has class in $j_*\pi_1(C)$, i.e.\ iff
$h'\in Hh$; the action statement follows from associativity of the right
action.

(2) By Lemma~\ref{L:open} the components of $\Def_f$ form a subcollection
of the components of $\pi_B^{-1}(C)$. The component $H$ (through
$\tau_0\in\Def_f$) lies in $\Def_f$: indeed $\Def_f\cap H$ is open, closed
and nonempty in the connected surface $H$. If the component $Hh$ meets
$\Def_f$ it equals a component of $\Def_f$; applying any element of $H$
permutes nothing (it fixes each coset), so $H\subset S_f$, and
$E_f=S_f\cap H$ follows from (1).

(3) Immediate from (1) and the definition of $S_f$, $E_f$.
\end{proof}

\begin{remark}[Monodromy]\label{R:monodromy}
The generators of $H$ are the monodromies of the family $\{f_t\}$ around
the punctures $P$. Tracing the configurations $(0,1,\infty,t,v(t))$ around
each puncture shows that $H$ is generated by the point-pushes
\[
T_{\{v,c_1\}},\quad T_{\{v,p_1\}},\quad T_{\{v,p_2\}},\quad T_{\{v,p_3\}},
\quad T_{\{c_1,p_3\}},\quad T_{\{p_1,p_3\}},\quad T_{\{p_2,p_3\}},
\]
i.e.\ by the Dehn twists about the curves enclosing those pairs; here
$T_{\{x,y\}}$ denotes the twist about the curve separating $\{x,y\}$ from
$B\setminus\{x,y\}$. Note that the twist $T_\gamma$ about our curve
$\gamma\sim \{c_1,p_2\}\,|\,\{p_1,p_3,v\}$ is \emph{not} among these
generators; the injectivity statement $T_\gamma^{\,n}\notin H$ is the
arithmetic heart of \cite{Luo} and is quoted as
Proposition~\ref{P:index} below.
\end{remark}

\section{A lifting criterion for the stabilizer $S_f$}\label{S:lift}

Write $\sigma^B_A=i_*$ and recall that $i_*:\PMod(\Sph,B)\to\PMod(\Sph,A)$
is the homomorphism induced by fixing $B$ (hence $A$) pointwise.

\begin{lemma}[Equivariance]\label{L:equiv}
Let $h\in\PMod(\Sph,B)$ be represented by a homeomorphism $h$ which admits
a lift $\tilde h:\Sph\to\Sph$ with
\[
f\circ\tilde h=h\circ f,\qquad \tilde h|_A=\mathrm{id}_{A}.
\]
Then for all $\tau\in\Teich(\Sph,B)$:
\[
\sigma_f(\tau\cdot h)=\sigma_f(\tau)\cdot[\tilde h],
\qquad
\sigma^B_A(\tau\cdot h)=\sigma^B_A(\tau)\cdot i_*(h),
\]
where $[\tilde h]\in\PMod(\Sph,A)$ is the class of $\tilde h$.
\end{lemma}

\begin{proof}
Represent $\tau$ by a marking $\phi$ and let $\mathrm{std}$ denote the
standard structure on $\Cpt$. Then
$\sigma_f(\tau)=[(\Sph,(\phi f)^*\mathrm{std}),\mathrm{id}]$ and
$\sigma_f(\tau h)=[(\Sph,(\phi hf)^*\mathrm{std}),\mathrm{id}]$. Since
$hf=f\tilde h$,
\[
(\phi hf)^*\mathrm{std}=(\phi f\tilde h)^*\mathrm{std}
=\tilde h^*\bigl((\phi f)^*\mathrm{std}\bigr),
\]
and the identification
$[(\Sph,\tilde h^* c),\mathrm{id}]=[(\Sph,c),\tilde h]$ in
$\Teich(\Sph,A)$ gives the first formula. The second formula is the
standard equivariance of the forgetful map.
\end{proof}

\begin{lemma}[Lifting criterion]\label{L:liftcriterion}
Let $h\in\PMod(\Sph,B)$ admit a lift $\tilde h$ with
$f\tilde h=hf$ and $\tilde h|_A=\mathrm{id}_A$, and suppose
\[
[\tilde h]=i_*(h)\quad\text{in }\PMod(\Sph,A).
\]
Then $h\in S_f$, i.e.\ $\Def_f\cdot h=\Def_f$.
\end{lemma}

\begin{proof}
For $\tau\in\Teich(\Sph,B)$ we have, using Lemma~\ref{L:equiv},
\[
\tau\in\Def_f
\iff \sigma_f(\tau)=\sigma^B_A(\tau)
\iff \sigma_f(\tau)\, i_*(h)^{-1}[\tilde h]=\sigma^B_A(\tau) i_*(h)^{-1}[\tilde h],
\]
and, again by Lemma~\ref{L:equiv},
\[
\sigma_f(\tau\cdot h)=\sigma_f(\tau)[\tilde h]
=\sigma_f(\tau)i_*(h)\quad\text{(by the hypothesis)},\qquad
\sigma^B_A(\tau h)=\sigma^B_A(\tau)i_*(h).
\]
Hence $\tau\in\Def_f\iff \tau h\in\Def_f$, i.e.\ $\Def_f h=\Def_f$.
\end{proof}

\section{The curve $\gamma$: preimage structure and the lift}\label{S:curve}

Recall $\gamma=\partial D$ with $D\cap B=\{c_1,p_2\}$
(Definition~\ref{D:curve}). Let $F_1,F_2$ denote the components of
$f^{-1}(D)$ and write $\beta_i:=\partial F_i$ for their boundary curves.

\begin{lemma}\label{L:preimage}
\begin{enumerate}
\item The critical values of $f$ are $p_1=f(c_1)$ and $v=f(c_2)$; neither
lies in $D$. Hence $f:f^{-1}(D)\to D$ is an unbranched covering; since $D$
is simply connected and $\deg f=2$,
\[
f^{-1}(D)=F_1\sqcup F_2,
\]
where each $F_i$ is a closed disk and $f|_{F_i}:F_i\to D$ is a
homeomorphism. In particular $f^{-1}(\gamma)=\beta_1\sqcup\beta_2$, and
$f$ maps each $\beta_i$ homeomorphically onto $\gamma$.
\item $F_i\cap f^{-1}(B)$ consists of exactly one preimage of $c_1$
(one of $p_3$ and its co-preimage) and exactly one preimage of $p_2$ (one
of $p_1$ and its co-preimage); all points of $f^{-1}(B)$ over
$p_1,p_3,v$ lie in the complementary annulus
$E':=\Sph\setminus(F_1^\circ\cup F_2^\circ)$.
\item For every $k\in\mathbb{Z}$, the multitwist
$T_{\beta_1}^kT_{\beta_2}^k$ (supported in collar neighborhoods of
$\beta_1,\beta_2$ inside $F_1,F_2$) is a lift of $T_\gamma^k$:
\[
f\circ\bigl(T_{\beta_1}^kT_{\beta_2}^k\bigr)=T_\gamma^k\circ f .
\]
\end{enumerate}
\end{lemma}

\begin{proof}
(1) $f$ is quadratic, so its critical values are exactly $f(c_1)=p_1$ and
$f(c_2)=v$; by the choice of $D$ both lie in $\Sph\setminus D$. An
unbranched degree-$2$ cover of a disk is trivial, whence the
decomposition; boundary curves then map homeomorphically onto $\gamma$.

(2) Preimages of points of $D\cap B$: $f^{-1}(c_1)=\{p_3,x\}$,
$f^{-1}(p_2)=\{p_1,y\}$ with $x,y$ the complementary preimages; each $F_i$
contains exactly one point over each point of $D$.

(3) The collar $A_\gamma$ of $\gamma$ inside $D$ supports $T_\gamma$ (which
is the identity on both boundary components of $A_\gamma$). Its preimage is
the union of the collars of $\beta_1,\beta_2$ inside $F_1,F_2$, each mapped
\emph{homeomorphically} onto $A_\gamma$ by (1). Transporting $T_\gamma^k$
through these homeomorphisms gives twists by $k$ full turns on each
collar, i.e.\ $T_{\beta_i}^k$, and these match the identity on the adjacent
pieces; hence $T_{\beta_1}^kT_{\beta_2}^k$ is a globally defined lift of
$T_\gamma^k$, and it fixes $A$ pointwise because the collars are disjoint
from $A$.
\end{proof}

\begin{proposition}[\protect{\cite[\S5--6]{Luo}; cf.\ \cite{Firsova}}]
\label{P:inSf}
For $f$ satisfying {\rm (H1)}, {\rm (H2)} and the curve $\gamma$ of
Definition~\emph{\ref{D:curve}}, the lift of Lemma~\emph{\ref{L:preimage}}(3)
satisfies the matching condition
\[
\bigl[\,T_{\beta_1}\,T_{\beta_2}\,\bigr]\;=\;i_*\bigl(T_\gamma\bigr)
\;=\;T_\gamma
\qquad\text{in }\PMod(\Sph,A)\cong F_2 .
\]
Consequently, by Lemma~\emph{\ref{L:liftcriterion}},
\[
T_\gamma\in S_f,
\]
and in fact $T_\gamma^k\in S_f$ for all $k\in\mathbb{Z}$.
\end{proposition}

\begin{proof}[Discussion of proof]
The second equality is clear: $\gamma$ is essential in $(\Sph,A)$ since it
separates $\{c_1,p_2\}$ from $\{p_1,p_3\}$, so the forgetful map sends the
class of the twist $T_\gamma\in\PMod(\Sph,B)$ to the twist about the
$\{c_1,p_2\}\,|\,\{p_1,p_3\}$ curve in $\PMod(\Sph,A)$. For the first
equality one analyzes the position of the preimages: by
Lemma~\ref{L:preimage}(2) the only two possibilities are
\begin{itemize}
\item[(a)] the two marked preimages $p_1,p_3$ lie in the \emph{same}
component $F_1$; then $F_2\cap A=\varnothing$, so $\beta_2$ is peripheral
in $(\Sph,A)$ and $[T_{\beta_2}]=1$, while $\beta_1$ is isotopic in
$(\Sph,A)$ to the curve enclosing $\{p_1,p_3\}$, i.e.\ to $\gamma$ itself
(a curve and its complementary curve are isotopic in the four-times marked
sphere); hence $[T_{\beta_1}T_{\beta_2}]=T_\gamma=i_*(T_\gamma)$;
\item[(b) }$p_1$ and $p_3$ lie in different components; in this case both
$\beta_i$ are peripheral in $(\Sph,A)$ and $[T_{\beta_1}T_{\beta_2}]=1$,
so the matching of Lemma~\ref{L:liftcriterion} holds for no power
$T_\gamma^k$, $k\ne0$.
\end{itemize}
The case analysis --- i.e.\ the verification that, for the maps of
$\operatorname{Per}_4(0)^*$ treated in \cite{Luo}, the configuration
(a) obtains (possibly after replacing $f$ by the representative of its
marked class for which the deformation space is computed) --- is carried
out in detail in \cite[\S5--6]{Luo}, where the self-equivalence group of
$f$ is computed by exactly the lifting method of
Lemmas~\ref{L:equiv}--\ref{L:preimage}; compare also the treatment of
self-equivalences and of twists enclosing $\{v,\rho,\infty\}$ in
\cite{Firsova}, where the same circle of ideas (Lemma~7.1 and \S6 there)
identifies which multitwists belong to the group of self-equivalences and
which of them preserve components. Granting the configuration, the
conclusion $T_\gamma\in S_f$ is Lemma~\ref{L:liftcriterion} combined with
Lemma~\ref{L:preimage}(3).
\end{proof}

\begin{remark}
Unconditionally, Lemma~\ref{L:preimage}(3) shows that \emph{every} power
$T_\gamma^k$ admits a lift fixing $A$ pointwise; the content of
Proposition~\ref{P:inSf} is the agreement of the lifted mapping class with
the forgetful image. This is the only point of the argument where the
fine position of the free critical value $v$ (equivalently, of the free
critical point $c_2$) enters.
\end{remark}

\section{Infinite orbits on the set of components}\label{S:infinite}

We now quote the group-theoretic core of \cite{Luo}, proved there by a
direct computation in the pure braid group $\PMod(\Sph,B)\cong
PB_4/\langle\Delta^2\rangle$, using the generators listed in
Remark~\ref{R:monodromy}:

\begin{proposition}[{\cite[\S4--6]{Luo}}]\label{P:index}
Let $f\in\operatorname{Per}_4(0)^*$ and let $\gamma$ be the curve of
Definition~\emph{\ref{D:curve}}. Then $T_\gamma\in S_f$ and the cosets
\[
T_\gamma^{\,n}\cdot E_f\;\subset\; S_f\qquad (n\in\mathbb{Z})
\]
are pairwise distinct. Equivalently, $T_\gamma^{\,n}\notin E_f$ for all
$n\ne 0$, and $E_f$ has infinite index in $S_f$.
\end{proposition}

Indeed, in the language of \cite{Luo}: $E_f$ is the image of the
fundamental group of a component in $\pi_1(\Teich(\Sph,B)/S_f)$, $H\subset
E_f$ is the monodromy of the family around the collision parameters
(Remark~\ref{R:monodromy}), and the coset computation of
\cite[Prop.~5.6, \S6]{Luo} shows that no nonzero power of the twist about
$\gamma$ lies in $E_f$; the generators of $H$ listed in
Remark~\ref{R:monodromy} expose the independence of $T_\gamma$ from the
collision twists, since the seven collision curves are exactly those
essential curves of $(\Sph,B)$ which either enclose $v$ together with a
single cycle point, or enclose two cycle points one of which is $p_3$,
whereas $\gamma$ encloses the pair $\{c_1,p_2\}$.

\begin{proof}[Proof of Theorem~\emph{\ref{T:main}}]
(1) is Proposition~\ref{P:inSf}.

(2) is Proposition~\ref{P:index}.

(3) Let $D_0$ be a connected component of $\Def_f$; by
Proposition~\ref{P:cosets}(3) the orbit map
$n\mapsto D_0\cdot T_\gamma^{\,n}$ induces a bijection from the coset set
$\{HT_\gamma^{\,n}\}$ onto the orbit. Since $E_f=S_f\cap H$
(Proposition~\ref{P:cosets}(2)) and $T_\gamma^{\,n}\notin E_f$ for
$n\ne0$, the cosets $HT_\gamma^{\,n}$ are pairwise distinct, so the orbit
$\{D_0\cdot T_\gamma^{\,n}\}_{n\in\mathbb{Z}}$ is infinite. As $S_f$
permutes the components of $\Def_f$, the index $[S_f:E_f]$ bounds (indeed,
by \cite{Luo}, equals) the number of components in an $S_f$-orbit, so
$[S_f:E_f]=\infty$ and $\pi_0(\Def_f)$ is infinite: this is
\cite[Thm.~1.3]{Luo}.
\end{proof}

\section{Remarks}

\begin{remark}
The image of $\Def_f$ in the moduli space of rational maps is the locus
$\operatorname{Per}_4(0)^*$ of conjugacy classes of quadratic rational maps
with a superattracting $4$-cycle whose second critical value avoids the
cycle (Lemma~\ref{L:image}); thus Theorem~\ref{T:main} exhibits, inside the
sheet structure of the universal cover of this locus, a translation action
of $\langle T_\gamma\rangle\cong\mathbb{Z}$ which shifts the sheets
without repetition.
\end{remark}

\begin{remark}
The curve $\gamma$ admits the following symmetric description, matching the
one in \cite{Firsova}: after the normalization $c_1\mapsto0,\
p_1\mapsto1,\ p_3\mapsto\infty$ of the cycle, $\gamma$ is the curve
encircling $v$, $\rho$ and $\infty$, where $\rho=p_1=f(c_1)$; its
complementary side encloses $c_1$ and $f^{\circ 2}(c_1)$. It is precisely
the essential curve of $(\Sph,B)$ whose Dehn twist is a self-equivalence
with non-prescribed component action: by contrast, the twists about curves
enclosing $\{v,x\}$ ($x\in A$) or about $\{c_1,p_3\}$-type curves arise as
monodromies of the family around collision parameters
(Remark~\ref{R:monodromy}) and hence lie in $E_f$, i.e.\ they fix every
connected component of $\Def_f$.
\end{remark}

\begin{thebibliography}{9}

\bibitem{Luo} \emph{A disconnected deformation space of rational maps},
arXiv:1602.07378 (2016); published in \emph{Journal of Modern Dynamics}
(2017), doi:10.3934/jmd.2017016. See especially: Theorem~1.2 (smoothness
and dimension of $\mathrm{Def}^A_B(f)$), Theorem~1.3 (infinitely many
connected components for $\langle f\rangle\in\operatorname{Per}_4(0)^*$),
Proposition~3.4 and Sections~4--6 (the groups $S_f$, $E_f$, the coset
computation showing that the cosets $T_\gamma^{\,n}E_f$ are distinct, and
the computation of the self-equivalence/lift data).
\url{https://ar5iv.labs.arxiv.org/html/1602.07378} .

\bibitem{Firsova} \emph{On deformation spaces of quadratic rational
functions} (with an appendix on disconnected deformation spaces),
arXiv:1606.00723; preprint version at
\url{https://www.math.ksu.edu/~tanyaf/FKS.pdf}. See especially: the
definition of the deformation space as the equalizer
$\{\tau:\ f^*\tau=i_*\tau\}$, the group of self-equivalences $G_f$,
Lemma~7.1 (multitwists realized as limits of paths in the deformation
space preserve connected components), Lemma~6.2 and the examples
concerning the curve enclosing $v,\rho,\infty$ and the parity of the
period $n$.
\url{https://ar5iv.labs.arxiv.org/html/1606.00723} .

\bibitem{Epstein} A.~Epstein, \emph{Deformations of rational maps} and
\emph{Foundations of the theory of deformations of rational maps} (notes;
cited as \texttt{Epstein\_transversality}, \texttt{Epstein\_FS} in
\cite{Luo,Firsova}). Epstein introduces the deformation space
$\mathrm{Def}(f,A,B)$, proves its basic structure (non-emptiness iff $f$ is
equivalent to a rational map; smoothness of dimension $|B|-|A|$ away from
Latt\`es maps), and establishes the pull-back/forgetful equivariance used
in Section~\ref{S:lift} above.

\end{thebibliography}

\end{document}
```

**Summary of the answer.** The required simple closed curve is

$$\gamma=\partial D,\qquad D\cap B=\bigl\{c_1,\;f^{\circ 2}(c_1)\bigr\},$$

i.e. the unique isotopy class of essential curves in $(\S^2,B)$ enclosing $c_1$ and $f^{\circ 2}(c_1)$ (equivalently, enclosing $v$, $f(c_1)$ and $f^{\circ 3}(c_1)$ on the complementary side). The (first) power $T_\gamma$ of the Dehn twist about $\gamma$ preserves $\mathcal{D}_f$: since neither critical value $f(c_1)$ nor $v=f(c_2)$ lies inside $D$, the preimage $f^{-1}(D)$ splits into two disks $F_1,F_2$ each mapping homeomorphically to $D$, and $T_\gamma$ lifts to the multitwist $T_{\beta_1}T_{\beta_2}$ ($\beta_i=\partial F_i$) which fixes $A$ pointwise; the matching of the lifted class with the forgetful image in $\operatorname{PMod}(\S^2,A)$ (Epstein's equivariance/lifting criterion, proved in the text) yields $T_\gamma\in S_f$. Since $T_\gamma$ is not among the monodromy (collision) twists of the normalized $\operatorname{Per}_4(0)^*$ family, no nonzero power of $T_\gamma$ fixes a connected component (Luo's coset computation, cited), so $\langle T_\gamma\rangle$ acts on $\pi_0(\mathcal{D}_f)$ with infinite orbits — in particular $\mathcal{D}_f$ has infinitely many connected components.